\section{Where the Research Can Be Used}

\textbf{Options trading.} Almost all stocks in the sample have NSE stock options. A one-day RV forecast
is the realized-volatility leg of the volatility risk premium: comparing it with the option's implied
volatility indicates whether straddles or strangles are rich or cheap, and it sets the expected gamma
P\&L of a delta-hedged position. The liquidity result matters here: for illiquid names the LSTM's edge is
largest, and these are also the options where implied volatility is least efficient.

\textbf{Risk management and margins.} Next-day volatility drives Value-at-Risk, Expected Shortfall and
exposure limits. Brokers and clearing members can use stock-level forecasts to set margins and position
limits; a forecast 6\% above realized on average (LSTM) is safer than one 21\% above (HAR-RV), which ties up
capital, or one that under-reacts.

\textbf{Position sizing and volatility targeting.} Scaling positions by forecast volatility improves
risk-adjusted returns \citep{fleming2001,moreira2017}. Better forecasts on ordinary days, where this
model gains most, translate directly into steadier risk.

\textbf{Execution and market making.} Expected intraday volatility sets quoted spreads, inventory limits
and how fast to work an order. The gain is largest in less liquid stocks, where execution cost is highest.

\textbf{Surveillance.} The gap between realized and forecast RV flags unusual days (for example, the Adani
episode in February 2023) for exchanges, regulators and risk desks.

\textbf{Model choice in practice.} For the 30 or so most liquid stocks, log-HAR is as good as the LSTM and
far simpler. A pooled LSTM pays for itself across a broad, mixed-liquidity universe. Because the model
loses ground on spike days, it should be paired with a stress overlay rather than used alone in crises.

\section{Limitations and Further Work}

\begin{itemize}
\item \textbf{One test period}, mostly calm, with no ``turbulent'' days by training standards. Results in a
crisis may differ; re-training on a rolling window would test this.
\item \textbf{Survivorship}: the universe is today's NIFTY~100, so stocks that dropped out are missing.
\item \textbf{Intraday RV only}: the overnight gap is excluded from the target. A close-to-close
volatility forecast would need an overnight component.
\item \textbf{Inputs}: the LSTM sees only RV history. India VIX, jump and semivariance measures
\citep{bollerslev2016}, volume and market-wide RV are natural additions.
\item \textbf{Testing}: a Model Confidence Set \citep{hansen2011} and an economic test (option P\&L or VaR
back-test) would complete the evaluation.
\end{itemize}

\section{Conclusion}

Across 95 NSE stocks and 15 months of out-of-sample data, an LSTM fed with HAR inputs forecasts next-day
realized volatility best: 15\% better than HAR-RV and 1.6\% better than log-HAR, both significant. The
network's own gain is modest and concentrated in less liquid stocks and calm-to-normal markets. Models
built only on daily returns (EGARCH-LSTM) are clearly worse, and EGARCH information adds nothing once the
intraday history is known. Careful data cleaning and leak-free design changed the size of the results
more than the choice of model did.
